% ======================================================================
% Beamer presentation - generated by KherveSlide.
% Edit freely; the source is regenerated when you change a slide.
% ======================================================================
\documentclass[aspectratio=169]{beamer}

% --- theme & packages ---
\usetheme{default}
\definecolor{ksth0}{HTML}{0E7D7D}
\setbeamercolor{structure}{fg=ksth0}
\definecolor{ksth1}{HTML}{212121}
\setbeamercolor{normal text}{fg=ksth1}
\definecolor{ksth2}{HTML}{0E7D7D}
\setbeamercolor{frametitle}{fg=ksth2}
\definecolor{ksth3}{HTML}{0E7D7D}
\setbeamercolor{title}{fg=ksth3}
\definecolor{ksth4}{HTML}{D4E8E8}
\setbeamercolor{block title}{bg=ksth4}
\definecolor{ksRule}{HTML}{85B09A}
\setbeamertemplate{footline}{%
\hbox{\color{ksRule}\rule{\paperwidth}{1.5pt}}\vskip0pt}
\usepackage{lmodern}
\usepackage{tikz}
\usetikzlibrary{arrows.meta}
\usetikzlibrary{shadows}
\usepackage[absolute,overlay]{textpos}
\usepackage{graphicx}
\setlength{\TPHorizModule}{1\paperwidth}
\setlength{\TPVertModule}{1\paperheight}
\textblockorigin{0\paperwidth}{0\paperheight}
\setbeamerfont{itemize/enumerate subbody}{size=\relax}
\setbeamerfont{itemize/enumerate subsubbody}{size=\relax}
\setbeamertemplate{footline}{%
\leavevmode\hbox{%
\begin{beamercolorbox}[wd=.333\paperwidth,ht=2.5ex,dp=1.2ex,leftskip=1.5ex]{author in head/foot}\end{beamercolorbox}%
\begin{beamercolorbox}[wd=.334\paperwidth,ht=2.5ex,dp=1.2ex,center]{title in head/foot}\end{beamercolorbox}%
\begin{beamercolorbox}[wd=.333\paperwidth,ht=2.5ex,dp=1.2ex,rightskip=1.5ex]{date in head/foot}\hfill \insertframenumber\end{beamercolorbox}}%
\vskip0pt}
\title{From molecule to model}
\author{Lab meeting}

\begin{document}

% ======================================================================
% Slide 1
% ======================================================================
\begin{frame}[t]

  % picture (free-positioned)
\begin{textblock}{0.14}(0.2,0.14)
\includegraphics[width=0.14\paperwidth,height=0.2489\paperheight,keepaspectratio]{media/icon_khervemol.png}
\end{textblock}

  % picture (free-positioned)
\begin{textblock}{0.14}(0.66,0.14)
\includegraphics[width=0.14\paperwidth,height=0.2489\paperheight,keepaspectratio]{media/icon_khervecad.png}
\end{textblock}

  % line / arrow (free-positioned)
\definecolor{ksline0}{HTML}{999999}
\begin{textblock}{0.24}(0.38,0.27)
\begin{tikzpicture}
\useasboundingbox (0,0) rectangle (\linewidth,-0\TPVertModule);
\draw[line width=3pt,color=ksline0, -{Stealth[length=2.4mm]}] (0\linewidth,-0\TPVertModule) -- (1\linewidth,-0\TPVertModule);
\end{tikzpicture}
\end{textblock}

  % text box (free-positioned)
\begin{textblock}{0.88}(0.06,0.52)
{\centering\fontsize{34}{41}\selectfont \textcolor[HTML]{00787A}{\textbf{From molecule to model}}\par}
\end{textblock}

  % text box (free-positioned)
\begin{textblock}{0.8}(0.1,0.68)
{\centering\fontsize{17}{20}\selectfont \textcolor[HTML]{555555}{Designing with KherveMol and KherveCAD}\par}
\end{textblock}
\end{frame}

% ======================================================================
% Slide 2 - Draw it: KherveMol
% ======================================================================
\begin{frame}[t]
  \frametitle{Draw it: KherveMol}

  % picture (free-positioned)
\begin{textblock}{0.6}(0.03,0.19)
\definecolor{ksBorder}{HTML}{BBBBBB}\begin{tikzpicture}
\node[inner sep=0pt,outer sep=0pt,draw=ksBorder,line width=0.8pt,drop shadow] (ksframe) {\includegraphics[width=0.6\paperwidth,height=0.62\paperheight,keepaspectratio]{media/mol_caffeine.jpg}};
\pgfresetboundingbox
\path[use as bounding box] (ksframe.south west) rectangle (ksframe.north east);
\end{tikzpicture}
\end{textblock}

  % text box (free-positioned)
\begin{textblock}{0.33}(0.65,0.22)
{\raggedright\fontsize{13}{16}\selectfont \begin{itemize}
\item Sketch in 2D, see it in 3D
\item SMILES and .mol / .sdf / .pdb import (RDKit)
\item Formula, mass, LogP, TPSA, InChI
\end{itemize}\par}
\end{textblock}

  % text box (free-positioned)
\begin{textblock}{0.6}(0.03,0.87)
{\centering\fontsize{11}{13}\selectfont \textcolor[HTML]{555555}{\textit{Caffeine in the 3D viewer}}\par}
\end{textblock}
\end{frame}

% ======================================================================
% Slide 3 - Relax it: geometry optimisation
% ======================================================================
\begin{frame}[t]
  \frametitle{Relax it: geometry optimisation}

  % picture (free-positioned)
\begin{textblock}{0.5}(0.04,0.2)
\includegraphics[width=0.5\paperwidth,height=0.62\paperheight,keepaspectratio]{media/example_energy.png}
\end{textblock}

  % text box (free-positioned)
\begin{textblock}{0.42}(0.55,0.26)
{\centering\fontsize{18}{22}\selectfont \[\mathbf{x}_{k+1} = \mathbf{x}_k - \mathbf{H}_k^{-1}\nabla E(\mathbf{x}_k)\]\par}
\end{textblock}

  % text box (free-positioned)
\begin{textblock}{0.38}(0.57,0.48)
{\begin{exampleblock}{Result}\raggedright\fontsize{13}{16}\selectfont Converged in 20 steps.\end{exampleblock}\par}
\end{textblock}
\end{frame}

% ======================================================================
% Slide 4 - Crystals, surfaces, nanostructures
% ======================================================================
\begin{frame}[t]
  \frametitle{Crystals, surfaces, nanostructures}

  % picture (free-positioned)
\begin{textblock}{0.42}(0.06,0.18)
\definecolor{ksBorder}{HTML}{BBBBBB}\begin{tikzpicture}
\node[inner sep=0pt,outer sep=0pt,draw=ksBorder,line width=0.8pt,drop shadow] (ksframe) {\includegraphics[width=0.42\paperwidth,height=0.28\paperheight,keepaspectratio]{media/mol_srtio3.jpg}};
\pgfresetboundingbox
\path[use as bounding box] (ksframe.south west) rectangle (ksframe.north east);
\end{tikzpicture}
\end{textblock}

  % text box (free-positioned)
\begin{textblock}{0.42}(0.06,0.515)
{\centering\fontsize{11}{13}\selectfont \textcolor[HTML]{555555}{\textit{SrTiO$_3$ perovskite}}\par}
\end{textblock}

  % picture (free-positioned)
\begin{textblock}{0.42}(0.52,0.18)
\definecolor{ksBorder}{HTML}{BBBBBB}\begin{tikzpicture}
\node[inner sep=0pt,outer sep=0pt,draw=ksBorder,line width=0.8pt,drop shadow] (ksframe) {\includegraphics[width=0.42\paperwidth,height=0.28\paperheight,keepaspectratio]{media/mol_mos2.jpg}};
\pgfresetboundingbox
\path[use as bounding box] (ksframe.south west) rectangle (ksframe.north east);
\end{tikzpicture}
\end{textblock}

  % text box (free-positioned)
\begin{textblock}{0.42}(0.52,0.515)
{\centering\fontsize{11}{13}\selectfont \textcolor[HTML]{555555}{\textit{MoS$_2$ (001) slab}}\par}
\end{textblock}

  % picture (free-positioned)
\begin{textblock}{0.42}(0.06,0.57)
\definecolor{ksBorder}{HTML}{BBBBBB}\begin{tikzpicture}
\node[inner sep=0pt,outer sep=0pt,draw=ksBorder,line width=0.8pt,drop shadow] (ksframe) {\includegraphics[width=0.42\paperwidth,height=0.28\paperheight,keepaspectratio]{media/mol_tbg.jpg}};
\pgfresetboundingbox
\path[use as bounding box] (ksframe.south west) rectangle (ksframe.north east);
\end{tikzpicture}
\end{textblock}

  % text box (free-positioned)
\begin{textblock}{0.42}(0.06,0.905)
{\centering\fontsize{11}{13}\selectfont \textcolor[HTML]{555555}{\textit{Twisted bilayer graphene}}\par}
\end{textblock}

  % picture (free-positioned)
\begin{textblock}{0.42}(0.52,0.57)
\definecolor{ksBorder}{HTML}{BBBBBB}\begin{tikzpicture}
\node[inner sep=0pt,outer sep=0pt,draw=ksBorder,line width=0.8pt,drop shadow] (ksframe) {\includegraphics[width=0.42\paperwidth,height=0.28\paperheight,keepaspectratio]{media/mol_c60.jpg}};
\pgfresetboundingbox
\path[use as bounding box] (ksframe.south west) rectangle (ksframe.north east);
\end{tikzpicture}
\end{textblock}

  % text box (free-positioned)
\begin{textblock}{0.42}(0.52,0.905)
{\centering\fontsize{11}{13}\selectfont \textcolor[HTML]{555555}{\textit{C$_{60}$}}\par}
\end{textblock}
\end{frame}

% ======================================================================
% Slide 5 - Build the hardware: KherveCAD
% ======================================================================
\begin{frame}[t]
  \frametitle{Build the hardware: KherveCAD}

  % picture (free-positioned)
\begin{textblock}{0.6}(0.03,0.19)
\definecolor{ksBorder}{HTML}{BBBBBB}\begin{tikzpicture}
\node[inner sep=0pt,outer sep=0pt,draw=ksBorder,line width=0.8pt,drop shadow] (ksframe) {\includegraphics[width=0.6\paperwidth,height=0.62\paperheight,keepaspectratio]{media/cad_vibe.jpg}};
\pgfresetboundingbox
\path[use as bounding box] (ksframe.south west) rectangle (ksframe.north east);
\end{tikzpicture}
\end{textblock}

  % text box (free-positioned)
\begin{textblock}{0.33}(0.65,0.22)
{\raggedright\fontsize{13}{16}\selectfont \begin{itemize}
\item Describe the part; Claude builds it
\item Real OpenSCAD nodes you keep editing
\item Parametric vacuum-hardware library
\item Print-check before you print
\end{itemize}\par}
\end{textblock}
\end{frame}

% ======================================================================
% Slide 6 - Assemble and document
% ======================================================================
\begin{frame}[t]
  \frametitle{Assemble and document}

  % picture (free-positioned)
\begin{textblock}{0.46}(0.03,0.2)
\definecolor{ksBorder}{HTML}{BBBBBB}\begin{tikzpicture}
\node[inner sep=0pt,outer sep=0pt,draw=ksBorder,line width=0.8pt,drop shadow] (ksframe) {\includegraphics[width=0.46\paperwidth,height=0.48\paperheight,keepaspectratio]{media/cad_exploded.jpg}};
\pgfresetboundingbox
\path[use as bounding box] (ksframe.south west) rectangle (ksframe.north east);
\end{tikzpicture}
\end{textblock}

  % picture (free-positioned)
\begin{textblock}{0.46}(0.51,0.2)
\definecolor{ksBorder}{HTML}{BBBBBB}\begin{tikzpicture}
\node[inner sep=0pt,outer sep=0pt,draw=ksBorder,line width=0.8pt,drop shadow] (ksframe) {\includegraphics[width=0.46\paperwidth,height=0.48\paperheight,keepaspectratio]{media/cad_blueprint.jpg}};
\pgfresetboundingbox
\path[use as bounding box] (ksframe.south west) rectangle (ksframe.north east);
\end{tikzpicture}
\end{textblock}

  % text box (free-positioned)
\begin{textblock}{0.46}(0.03,0.74)
{\centering\fontsize{11}{13}\selectfont \textcolor[HTML]{555555}{\textit{Exploded view of a vacuum stack}}\par}
\end{textblock}

  % text box (free-positioned)
\begin{textblock}{0.46}(0.51,0.74)
{\centering\fontsize{11}{13}\selectfont \textcolor[HTML]{555555}{\textit{Automatic third-angle drawing}}\par}
\end{textblock}
\end{frame}

% ======================================================================
% Slide 7 - Into the lab
% ======================================================================
\begin{frame}[t]
  \frametitle{Into the lab}

  % picture (free-positioned)
\begin{textblock}{0.6}(0.03,0.19)
\definecolor{ksBorder}{HTML}{BBBBBB}\begin{tikzpicture}
\node[inner sep=0pt,outer sep=0pt,draw=ksBorder,line width=0.8pt,drop shadow] (ksframe) {\includegraphics[width=0.6\paperwidth,height=0.66\paperheight,keepaspectratio]{media/cad_bench.jpg}};
\pgfresetboundingbox
\path[use as bounding box] (ksframe.south west) rectangle (ksframe.north east);
\end{tikzpicture}
\end{textblock}

  % text box (free-positioned)
\begin{textblock}{0.31}(0.66,0.3)
{\begin{block}{Idea}\raggedright\fontsize{14}{17}\selectfont Glassware, benches and molecules share one scene — plan an experiment before you set it up.\end{block}\par}
\end{textblock}
\end{frame}

\end{document}
